\chapter{Alphabete, Wörter, Sprachen}
Wenn man sich mit der Funktionsweise von Rechnern (Computern) genauer beschäftigt, stellt man fest, dass Rechner im Grunde eine Transformation von Eingabedaten in Ausgabedaten realisieren. Sowohl die Eingabedaten als auch die Ausgabedaten lassen sich als Texte darstellen. Die Texte sind nichts anderes als Folgen von Symbolen aus einem bestimmten Alphabet. Programme können als Folge von Symbolen der Computertastatur dargestellt werden. In digitalen Rechnern sind alle Informationen als Folgen von Einsen und Nullen gespeichert. Damit realisiert der Rechner eine Transformation von Eingabetexten in Ausgabetexte.

In diesem Kapitel wollen wir den Formalismus für den Umgang mit Texten kennenlernen. Wir werden die fundamentalen Begriffe \textit{Alphabet}, \textit{Wort} und \textit{Sprache} einführen. Diese werden uns später helfen, bekannte Probleme der Informatik wie beispielsweise das \textit{Entscheidungsproblem} mathematisch sauber zu formulieren.

Dieses Kapitel folgt dem Abschnitt 2.2 aus dem Buch\footnote{J. Hromkovic: Theoretische Informatik. 5. Auflage, Springer Vieweg 2014., ISBN: 978-3-658-06432-7} sehr  nahe. Es wurden einige Aufgaben, welche als Hilfestellung dienen, hinzugefügt und kleine Teile ausgelassen.

\section{Alphabete, Wörter, Sprachen}
\begin{mydefinition}[Alphabet]\label{mydefinition:Alphabet}
Eine \textbf{endliche nichtleere} Menge $\Sigma$ heisst \textit{Alphabet}. Die \textit{Elemente eines Alphabets} werden \textit{Buchstaben} (\textit{Zeichen, Symbole}) genannt.
\end{mydefinition}
Wir werden später Alphabete verwenden, um eine schriftliche Darstellung einer Sprache zu erzeugen. \cref{mydefinition:Alphabet}, entspricht unserer intuitiven Vorstellung eines Alphabets:
Um Text darstellen zu können, muss ein Alphabet mindestens ein Symbol enthalten ($\to$ nichtleer). Damit man sich auf einen fixen Satz von Zeichen einigen kann, darf das Alphabet nicht unendlich gross sein ($\to$ endlich).

Wir listen nun einige der Alphabete auf, die in der Mathematik und Informatik häufig verwendet werden.
\begin{myexample}
\begin{aenum}
	\item $\Sigma_{\text{bool}} = \lrc{0,1}$ ist das Boole'sche Alphabet, mit dem digitale Rechner arbeiten.
	\item $\Sigma_{\text{lat}} = \lrc{a,b,c,\ldots ,z,A,B,\ldots ,Z}$ ist das lateinische Alphabet.
	\item $\Sigma_{\text{Tastatur}} = \lrc{a,b,c,\ldots ,z,A,B,C,\ldots ,Z,\text{\textvisiblespace}\,,>,<,(,),\ldots ,\#,?,!}$ ist das Alphabet aller Symbole, die mit der englischen Tastatur getippt werden können. Dabei ist \textvisiblespace\ das Symbol für das Leerzeichen / Leersymbol.
	\item $\Sigma_{\text{greek}} = \lrc{\alpha, \beta, \gamma, \ldots , \omega, A, B, \Gamma, \ldots , \Omega}$ ist das griechische Alphabet.
	\item $\Sigma_m = \setcm{n\in\N}{n<m}$ für jede fixe Wahl von $m\in\mathbb{N}\setminus\lrc{0}$, ist ein Alphabet für die $m$-adische Darstellung von Zahlen.
\end{aenum}
\end{myexample}

\textit{Wörter} werden wir als endliche Folgen von Buchstaben ansehen.
\begin{mydefinition}[Wort]\label{mydefinition:Wort}
Sei $\Sigma$ ein Alphabet. Ein \textit{Wort} über $\Sigma$ ist eine endliche (möglicherweise leere) Folge von Buchstaben aus $\Sigma$. Das \textit{leere Wort} $\lambda$ ist die leere Buchstabenfolge.
Die \textit{Länge} $|w|$ eines Wortes $w$ ist die Länge des Wortes als Folge, das heisst die Anzahl der Vorkommen von Buchstaben in $w$.
\end{mydefinition}

\begin{myexample}
\begin{aenum}
	\item $w = 1,0,0,1,0$ ist ein Wort über dem Alphabet $\Sigma_{\text{bool}}$ und $\abs{w} = 5$, da $w$ eine Folge von $5$ Buchstaben ist.
	\item $u = M,\text{\textvisiblespace}\,, j$ ist ein Wort über $\Sigma_{\text{Tastatur}}$ und $\abs{u} = 3$, da $u$ eine Folge von $3$ Buchstaben ist.
	\item Das leere Wort $\lambda$ ist ein Wort über jedem Alphabet und es gilt $\abs{\lambda} = 0$.
\end{aenum}
\end{myexample}
\begin{myexercise}\label{myexercise:0200}
	Was ist $\Sigma_{10}$?
\end{myexercise}

\begin{mydefinition}[Stern-Operator]
Sei $\Sigma$ ein Alphabet. $\Sigma^*$ (gesprochen: \textit{Sigma Stern}) ist \textbf{die Menge aller Wörter} über $\Sigma$, also die Menge aller endlichen Folgen von Symbolen aus $\Sigma$. Man nennt $\Sigma^{*}$ auch den \textbf{Kleene'schen Stern}\footnote{benannt nach dem US-Amerikanischen Mathematiker Stephen Cole Kleene} von $\Sigma$. $\Sigma^+ = \Sigma^*\setminus\lrc{\lambda}$ ist die Menge aller Wörter über $\Sigma$ ohne das leere Wort.
\end{mydefinition}

Wir werden im Folgenden Wörter ohne Kommmas schreiben. Anstelle von $1,0,0,1,0$ werden wir lediglich $10010$ schreiben. Allgemein, werden wir anstelle von $x_1,x_2,\ldots ,x_n$ einfach $x_1x_2\ldots x_n$ schreiben.
\begin{myremark}
	In der deutschen Sprache existieren die zwei Ausdrücke \enquote{Wörter} und \enquote{Worte}. Dies sind keine Synonyme. \textit{Wörter} bezeichnet den Plural von \textit{Wort} (\enquote{Im Duden stehen viele Wörter}). Der Ausdruck \textit{Worte} bezieht sich auf Gedankenkonstrukte (\enquote{Sie sprach weise Worte}).
\end{myremark}

Wir können Wörter benutzen, um mathematische Objekte wie Zahlen, Formeln, Graphen und Computerprogramme darzustellen. Ein Wort $x = x_1x_2\ldots x_n\in \lr{\Sigma_{\text{bool}}}^*$ kann als binäre Darstellung der Zahl
\begin{align}\label{eq:Nummer}
	\text{Nummer}(x)=\displaystyle\sum_{k=1}^n x_k\cdot 2^{n-k}
\end{align}
betrachtet werden.

\begin{myexercise}\label{myexercise:0298}
	Sei $x = 1011$. Berechnen Sie $\text{Nummer}(x)$.
\end{myexercise}
Für eine Zahl $m\in\N\setminus\lrc{0}$ wird mit $\text{Bin}(m)\in \lr{\Sigma_{\text{bool}}}^*$ die kürzeste binäre Darstellung von $m$ bezeichnet, also gilt $\text{Nummer}\lr{\text{Bin}(m)} = m$. Man setzt $\text{Bin}(0) = 0$.

% \begin{myexercise}%[difficulty={1}]
% 	Sei $m\in\mathbb{N}-\lrc{0}$. Welche Bedingung muss gelten, damit $\text{Bin}(m)$ die \textbf{kürzeste} binäre Darstellung von $m$ ist.\\
% \end{myexercise}
% \begin{myanswer}
% 	Die Bedingung ist lediglich, dass die Darstellung mit dem Symbol $1$ beginnt, da führende Nullen die Bedeutung der Darstellung nicht verändern, sondern die Darstellung lediglich länger machen. Zum Beispiel beschreiben $101$, $0101$ und $000000101$ dieselbe Zahl. Die erste ist aber die kürzeste Darstellung.
% \end{myanswer}
%
% \begin{theo}[Länge der binären Darstellung]{theo:laenge_binaer}
% 	Sei $m\in\mathbb{N}-\lrc{0}$. Dann ist die Länge $|\text{Bin}(m)|$ der kürzesten binären Darstellung von $m$ gegeben durch
% 	\begin{align*}
% 		|\text{Bin}(m)| = \ceil{\log_2(m+1)}
% 	\end{align*}
% \end{theo}
% Dabei bezeichnet für eine reelle Zahl $x$ der Ausdruck $\ceil{x}$ die nächste ganze Zahl $\geq x$. Genauer: Sei $x$ eine reelle Zahl, dann definiert man
% \begin{align*}
% 	\ceil{x} = \min\setcm{k\in\mathbb{Z}}{k\geq x}
% \end{align*}
% und analog
% \begin{align*}
% 	\floor{x} = \max\setcm{k\in\mathbb{Z}}{k\leq x}.
% \end{align*}
% \begin{myexercise}%[difficulty={1}]
% 	Bestimmen Sie mithilfe der Formel aus Satz \cref{theo:laenge_binaer} die Länge der kürzesten binären Darstellung der Zahl 58902.
% \end{myexercise}
% \begin{myanswer}
% 	Mithilfe eines Taschenrechners finden wir
% 	\begin{align*}
% 		\log_2(58902 + 1)\approx 15.846
% 	\end{align*}
% 	und damit
% 	\begin{align*}
% 		\ceil{\log_2(58902 + 1)} = 16
% 	\end{align*}
%
% \end{myanswer}
%
% \clearpage
% \begin{myexercise}%[difficulty={1}]
% 	Beweisen Sie Satz \cref{theo:laenge_binaer}.

% 	Tipp:
% 	\begin{itemize}
% 		\item Setzen Sie $k = |\text{Bin}(m)|$.
% 		\item Erklären Sie, warum die Sandwich Beziehung $2^k-1 \geq m > 2^{k-1} - 1$ gilt, indem Sie folgende Fragen beantworten:
% 		      \begin{itemize}
% 			      \item Was ist der Dezimalwert der kleinsten binären Zahl der Länge $k+1$ (ohne führende Nullen)?
% 			      \item Was ist der Dezimalwert der kleinsten binären Zahl der Länge $k$ (ohne führende Nullen)?
% 		      \end{itemize}
% 	\end{itemize}
% \end{myexercise}
% \begin{myanswer}
% 	Sei $k = |\text{Bin}(m)|$. Wir wollen $k$ ermitteln.
%
% 	Die kleinste binäre Zahl der Länge $k+1$ ist $1\underbrace{00\ldots 0}_{k\;\text{mal}}$. Diese hat gemäss der Gleichung \cref{eq:Nummer} den dezimalen Wert (Nummer) $2^{k}$. Dann kann der Dezimalwert der grössten binäre Zahl der Länge $k$ höchstens $2^{k-1}-1$ sein ($1$ weniger). Damit ist $2^k-1 \geq m$ begründet.
%
% 	Die kleinste binäre Zahl der Länge $k$ ist $1\underbrace{00\ldots 0}_{(k-1)\;\text{mal}}$, was dem dezimalen Wert $2^{k-1}$ entspricht. Also muss auch $m > 2^{k-1} - 1$ gelten, was die Sandwich-Beziehung aus dem Tipp begründet.\\
%
% 	Offensichtlich gilt nun also für $m$ die Beziehung
% 	\begin{align*}
% 		2^k-1 \geq m > 2^{k-1} - 1 \Rightarrow \\
% 		2^k \geq m+1 > 2^{k-1} \Rightarrow     \\
% 		k \geq \log_2(m+1) > k-1.
% 	\end{align*}
% 	Die Beziehung $k \geq \log_2(m+1) > k-1$ sagt aus, dass $\log_2(m+1)\in (k-1,k]$ und damit folgt:
% 	\begin{align*}
% 		k = \ceil{\log_2(m+1)}.
% 	\end{align*}
% \end{myanswer}

\begin{myexercise}\label{myexercise:0201}
	Wie sehen die folgenden Mengen aus?
	\begin{aenum}
		\item $\lrc{1}^{*}$
		\item $\lr{\Sigma_{\text{bool}}}^{*}$
	\end{aenum}
\end{myexercise}


\begin{myexercise}\label{myexercise:0202}
	In \cref{mydefinition:Wort}, haben wir ein Wort als endliche Folge von Buchstaben definiert. Nun enthält zum Beispiel die Menge $\lrc{1}^{*}$ in \cref{myexercise:0201} aber Wörter beliebiger Länge. Erklären Sie, warum dies kein Widerspruch ist.
\end{myexercise}

Da Alphabete insbesondere auch Mengen sind, kann man auch das kartesische Produkt von Alphabeten bilden.

\begin{myexample}
Sei $\Sigma_1 = \lrc{g,h}$ und $\Sigma_2 = \lrc{x,y}$, dann ist das kartesische Produkt $\Sigma_1\times \Sigma_2 = \lrc{(g,x),(g,y),(h,x),(h,y)}$.
\end{myexample}

Wir werden nun eine Operation einführen, welche uns erlaubt Wörter zu verketten. Diese Operation werden wir sehr häufig verwenden.
\begin{mydefinition}[Verkettung / Konkatenation]\label{mydefinition:Konkat}
Die \textbf{Verkettung} (\textbf{Konkatenation}) für ein Alphabet $\Sigma$ ist eine Abbildung
Kon: $\Sigma^*\times \Sigma^*\rightarrow\Sigma^*$, sodass
\begin{align*}
	\text{Kon}(x,y) = x\cdot y = xy
\end{align*}
für alle $x,y\in\Sigma^{*}$.
\end{mydefinition}
\begin{myexample}
Sei $\Sigma = {8,9,d,e}$ und seien $x = e899dd$ und $y = de8$, dann ist $\text{Kon}(x,y) = x\cdot y = e899ddde8$.
\end{myexample}

Wir werden fast ausnahmslos $xy$ schreiben und nur selten $x\cdot y$ oder $\text{Kon}(x,y)$.

\begin{myexercise}\label{myexercise:0203}
	Erläutern Sie in Ihren eigenen Worten, was die Bedeutung des Ausdrucks $\Sigma^*\times \Sigma^*\rightarrow\Sigma^*$ in \cref{mydefinition:Konkat} ist.
\end{myexercise}

Die Verkettung $\text{Kon}$ über $\Sigma$ ist assoziativ über $\Sigma^*$, da offensichtlich
\begin{align*}
	\text{Kon}(x,\text{Kon}(y,z)) = x\cdot (y\cdot z) = xyz = (x\cdot y)\cdot z = \text{Kon}(\text{Kon}(x,y),z)
\end{align*}
gilt, für alle $x,y,z\in\Sigma^*$.

\begin{myexercise}\label{myexercise:0204}
	Bei der Multiplikation in den reellen Zahlen ist die Zahl $1\in\R$ das neutrale Element, da $1\cdot x = x\cdot 1 = x$ für alle $x\in\R$ gilt.
	Welches ist das neutrale Element der Addition in den reellen Zahlen?
\end{myexercise}

Für jedes $x \in\Sigma^*$ gilt
\begin{align*}
	x\cdot\lambda = \lambda\cdot x = x.
\end{align*}
Damit ist $\lambda$ das neutrale Element der Verkettung über $\Sigma^*$.

\begin{myexercise}\label{myexercise:0205}
	Sei $x = a01b$ und $y = c21$. Bestimmen Sie $\abs{x}$, $\abs{y}$ und $\abs{xy}$.
	Seien allgemein $x,y\in\Sigma^*$ zwei Wörter für ein Alphabet $\Sigma$. Finden Sie einen Ausdruck für $\abs{xy}$.
\end{myexercise}


\begin{mydefinition}[Umkehrung / Reversal]\label{mydefinition:reversal}
Sei $n$ eine natürliche Zahl. Für ein Wort $x=x_1x_2\ldots x_n$, mit $x_i\in\Sigma$ für $i\in\lrc{1,2,\ldots ,n}$ bezeichnet $x^R = x_nx_{n-1}\ldots x_1$ die \textbf{Umkehrung} (\textbf{Reversal}) von $x$.
\end{mydefinition}

\begin{myexample}
Es sei $w := abcde$. Dann gilt $w^R = edcba$.
\end{myexample}

\begin{myexercise}\label{myexercise:0206}
	Sei $\Sigma$ ein Alphabet und $u,v\in\Sigma^*$ zwei Wörter. Beweisen oder widerlegen Sie die Aussage:
	\begin{align*}
		(uv)^R = v^Ru^R
	\end{align*}
\end{myexercise}

\begin{mydefinition}[Iteration eines Wortes]\label{mydefinition:Iteration_Wort}
Sei $\Sigma$ ein Alphabet. Für alle $x\in\Sigma^*$ und alle $i\in\mathbb{N}$ definieren wir die $i$-te \textbf{Iteration} $x^i$ von $x$ als
\begin{align*}
	x^0 = \lambda,\quad x^1 = x, \quad x^i = xx^{i-1}.
\end{align*}
\end{mydefinition}
\begin{myexample}
Sei $\Sigma = \lrc{a,b,c}$. Wir können nun schreiben:
\begin{align*}
	aa       & = a^2                 \\
	abababab & = (ab)^4              \\
	cbbbbbab & = cb^5ab = c(bb)^2bab
\end{align*}
\end{myexample}

Mit der folgenden Definition wollen wir den Begriff \textit{Teilwort}, für den wir eine gute Intuition haben, formalisieren. Ein Teilwort eines Wortes $x$ ist ein zusammenhängender Teil von $x$.
\begin{mydefinition}[Teilwörter]\label{mydefinition:Teilwort}
Seien $u,w\in\Sigma^*$ für ein Alphabet $\Sigma$.
\begin{itemize}
	\item $v$ heisst \textbf{Teilwort} von $w$ $\Leftrightarrow$ es existieren $x,y\in\Sigma^*$, sodass $w = xvy$
	\item $v$ heisst \textbf{Präfix} von $w$ $\Leftrightarrow$ es existiert $x\in\Sigma^*$, sodass $w = vx$
	\item $v$ heisst \textbf{Suffix} von $w$ $\Leftrightarrow$ es existiert $x\in\Sigma^*$, sodass $w = xv$
	\item $v\neq \lambda$ heisst \textbf{echtes} Teilwort (Präfix, Suffix) von $w$ $\Leftrightarrow$ $v\neq w$ und $v$ ein Teilwort (Präfix, Suffix) von $w$ ist
\end{itemize}
\end{mydefinition}

\begin{myexercise}\label{myexercise:0299}
	Es seien $\Sigma := \lrc{a,b,c}$ und $w := abc$ ein Wort über $\Sigma$. Bestimmen Sie alle Teilwörter von $w$. Welche dieser Teilwörter sind auch Präfixe von $w$?
\end{myexercise}

\begin{myexample}
Sei $\Sigma = \lrc{a,b,c}$. Das Wort $abc$ ist ein echtes Teilwort, echtes Präfix und ein echtes Suffix von $(abc)^3$. Das leere Wort $\lambda$ ist Teilwort von jedem Wort. Jedes Wort ist Teilwort von sich selbst (aber kein echtes Teilwort).
\end{myexample}

\begin{mychallenge}\label{mychallenge:0207}
	Es sei $x$ ein Wort der Länge $n\in\N$, welches aus lauter verschiedenen Buchstaben besteht (also aus $n$ verschiedenen Buchstaben). Wie viele verschiedene Teilwörter hat $x$?
\end{mychallenge}

Wir wollen an dieser Stelle noch eine weitere nützliche Schreibweise definieren.

Sei $x\in\Sigma^*$ und $a\in\Sigma$, dann ist $\abs{x}_a$ definiert als die Anzahl der Vorkommen von $a$ in $x$.
\begin{myexample}
Sei $v = babaac$, dann ist $\abs{v}_a = 3$, $\abs{v}_b = 2$ und $\abs{v}_c = 1$. Offensichtlich gilt für alle $x\in\Sigma^*$
\begin{align*}
	\abs{x} = \sum_{a\in\Sigma}\abs{x}_a.
\end{align*}
\end{myexample}

Nun kommen wir zur Definition einer Sprache. Dies wird für uns einer der wichtigsten Begriffe sein.
\begin{mydefinition}[Sprache]\label{mydefinition:Sprache}
Eine \textbf{Sprache} $L$ über einem Alphabet $\Sigma$ ist eine Teilmenge von $\Sigma^*$. Das Komplement $L^C$ der Sprache $L$ bezüglich $\Sigma$ ist die Sprache $\Sigma^*\setminus L$.
\begin{itemize}
	\item $L_{\emptyset} = \emptyset$ ist die leere Sprache.
	\item $L_{\lambda} = \lrc{\lambda}$ ist die einelementige Sprache, die nur aus dem leeren Wort besteht.
\end{itemize}
Sind $L_1$ und $L_2$ Sprachen über $\Sigma$, so bezeichnet
\begin{align*}
	L_1\cdot L_2 = L_1L_2 = \setcm{vw}{v\in L_1\; \text{und}\; w\in L_2}
\end{align*}
die \textbf{Konkatenation} von $L_1$ und $L_2$.

Ist $L$ eine Sprache über $\Sigma$, so definieren wir die Iterationen
\begin{align*}
	L^0 & = L_{\lambda}, L^{i+1} = L^i\cdot L\;\text{für alle}\; i\in\mathbb{N},                \\
	L^* & = \bigcup_{i\in\mathbb{N}}L^i\;\text{und}\; L^+ = \bigcup_{i\in\mathbb{N}}L^i\cdot L.
\end{align*}
\end{mydefinition}

\begin{myexample}
Die folgenden Mengen sind Beispiele von Sprachen über $\Sigma = \lrc{a,b}$.
\begin{itemize}
	\item $L_1 = \emptyset$
	\item $L_2 = \lrc{\lambda}$
	\item $L_3 = \Sigma^* = \lrc{\lambda, a,b,aa,\ldots }$
	\item $L_4 = \Sigma^+ = \lrc{a,b,aa,\ldots }$
	\item $L_5 = \Sigma$
	\item $L_6 = \setcm{a^p}{p\;\text{ist eine Primzahl}}$
	\item $L_5 = \lrc{a}^* = \lrc{\lambda, a, aa, aaa, aaaa, \ldots }$
	\item $\Sigma^2 = \lrc{aa,ab,ba,bb}$
\end{itemize}
\end{myexample}
Man beachte, dass $\Sigma^i = \setcm{x\in\Sigma^*}{\abs{x} = i}$ für eine natürliche Zahl $i$, und dass $L_{\emptyset}\cdot L = \emptyset$, $L_{\lambda}\cdot L = L$.

\begin{myexample}\leavevmode
\begin{itemize}
	\item Die Menge aller grammatikalisch korrekten englischen Texte ist eine Sprache über $\Sigma_{\text{Tastatur}}$.
	\item Die Menge aller syntaktisch korrekten Programme in {\tt C++} ist auch eine Sprache über $\Sigma_{\text{Tastatur}}$.
\end{itemize}
\end{myexample}

\begin{myexercise}\label{myexercise:0208}
	Sei $L_1 = \lrc{\lambda, a, b}$ und sei $L_2 = \lrc{a^4, a^2b}$. Welche Wörter Liegen in der Sprache $L_3 = L_1L_2$?
\end{myexercise}

\begin{mychallenge}\label{mychallenge:0209}
	Sei $k\in\mathbb{N}$. Geben Sie ein Alphabet $\Sigma$ und zwei Sprachen $L_1$ und $L_2$ über $\Sigma$ an, sodass
	\begin{align*}
		\abs{L_1} = k\quad \text{und} \quad \abs{L_1L_2} = k + 1.
	\end{align*}
\end{mychallenge}


\clearpage
\section{Lösungen der Aufgaben}
\begin{myanswer}[myexercise:0200]
	\begin{align*}
		\Sigma_{10} = \lrc{0,1,2,\ldots ,9}
	\end{align*}
	ist das Alphabet für die Darstellung von Zahlen im Dezimalsystem.
\end{myanswer}


\begin{myanswer}[myexercise:0298]
	\begin{align*}
		\text{Nummer}(x) = \text{Nummer}(\red{1}\blue{0}\green{1}\orange{1}) = \displaystyle\sum_{k=1}^4 x_k\cdot 2^{4-k} = \red{1}\cdot 2^3+\blue{0}\cdot 2^2 + \green{1}\cdot 2^1+\orange{1}\cdot 2^0 = 8 + 2 + 1 = 11
	\end{align*}
\end{myanswer}


\begin{myanswer}[myexercise:0201]
	\begin{align*}
		\lrc{1}^* & = \lrc{\lambda, 1, 11, 111, 1111, 11111, \ldots } =                                                       \\
		          & = \lrc{\lambda}\cup \setcm{x_1x_2\ldots  x_i}{i\in\mathbb{N},\; x_j = 1\;\text{für}\; j = 1,2,\ldots , i}
	\end{align*}
	\begin{align*}
		\lr{\Sigma_{\text{bool}}}^* = \lrc{0,1}^* & = \lrc{\lambda, 0, 1, 00, 01, 10, 11, 000, 001, 010, 100, 011, \ldots } =                                                     \\
		                                          & = \lrc{\lambda}\cup \setcm{x_1x_2\ldots  x_i}{i\in\mathbb{N},\; x_j\in\Sigma_{\text{bool}} \;\text{für}\; j = 1,2,\ldots , i}
	\end{align*}
\end{myanswer}

\begin{myanswer}[myexercise:0202]
	Anhand dieser Frage lässt sich der Unterschied zwischen den Begriffen \textit{unbeschränkt} und \textit{unendlich} schön verdeutlichen. Betrachten wir dazu exemplarisch nochmals die Menge $\lrc{1}^*$ aus \cref{myexercise:0201}:
	\begin{align*}
		\lrc{1}^* & = \lrc{\lambda, 1, 11, 111, 1111, 11111, \ldots } =                                                      \\
		          & = \lrc{\lambda}\cup \setcm{x_1x_2\ldots x_i}{i\in\mathbb{N},\; x_j = 1\;\text{für}\; j = 1,2,\ldots , i}
	\end{align*}
	\begin{itemize}
		\item Für jede natürliche Zahl $n\in\mathbb{N}$, existiert in $\lrc{1}^*$ ein Wort $w$ mit $\abs{w} \geq n$, da für jedes $n\in\mathbb{N}$ (insbesondere) auch das Wort $w:=1^n$ in der Menge $\lrc{1}^*$ enthalten ist und $\abs{w}\geq n$. Damit gibt es \textbf{keine obere Schranke} für die Länge der Wörter in $\lrc{1}^*$.
		\item Jedes Wort $w\in\lrc{1}^*$ hat die Form $w = 1^n$ für eine natürliche Zahl $n$. Damit hat jedes Wort in der Menge $\lrc{1}^*$ eine endliche Länge.
	\end{itemize}
	Die Länge der Wörter in $\lrc{1}^*$ können also nicht nach oben beschränkt werden, trotzdem hat jedes Wort in $\lrc{1}^*$ eine endliche Länge und ist somit tatsächlich ein Wort im Sinne von \cref{mydefinition:Wort}.
\end{myanswer}

\begin{myanswer}[myexercise:0203]
	Wir haben zwei Wörter $x$ und $y$ über dem Alphabet $\Sigma$. Weil $\Sigma^*$ die Menge aller Wörter über $\Sigma$ ist, gilt offensichtlich $x,y\in\Sigma^*$. Aus den zwei Wörtern $x$ und $y$ wird nun ein neues Wort $xy$ gebildet. Man beachte, dass die Reihenfolge $xy\neq yx$ (im Allgemeinen) eine Rolle spielt.	Die Verkettung nimmt also ein geordnetes Paar $(x,y)\in \lr{\Sigma^*\times \Sigma^*}$ (dies ist gerade die Menge aller geordneten Paare von Wörtern über $\Sigma$) und bildet dieses auf ein neues Wort $xy\in\Sigma^*$ ab.
\end{myanswer}

\begin{myanswer}[myexercise:0204]
	Die Zahl $0\in\R$, da $0+x = x+0 = x$ für alle $x\in\R$.
\end{myanswer}

\begin{myanswer}[myexercise:0205]
	Offensichtlich sind $\abs{x} = \abs{a01b}= 4$ und $\abs{y} = \abs{c21} = 3$ und somit $\abs{xy} = \abs{a01bc21} = 7$.

	Seien  $x,y\in\Sigma^*$ zwei Wörter über einem Alphabet $\Sigma$. Dann gilt
	\begin{align*}
		\abs{xy} = \abs{x}+\abs{y}.
	\end{align*}
\end{myanswer}

\begin{myanswer}[myexercise:0206]
	Da $u$ und $v$ Wörter sind, haben sie endliche Längen. Wir definieren $n$ und $m$ als $n := |u|$ und $m := |v|$. Dann hat $u$ die Form $u = u_1u_2\ldots u_n$ und $v = v_1v_2\ldots v_m$, wobei $u_1,u_2,\ldots ,u_n,v_1,v_2,\ldots ,v_m\in\Sigma$. Damit ist
	\begin{align*}
		uv = u_1u_2\ldots u_nv_1v_2\ldots v_m
	\end{align*}
	und somit
	\begin{align*}
		(uv)^R = v_m\ldots v_2v_1u_n\ldots u_2u_1 = v^Ru^R,
	\end{align*}
	da $v^R = v_m\ldots v_2v_1$ und $u^R = u_n\ldots u_2u_1$ gilt.
\end{myanswer}


\begin{myanswer}[myexercise:0299]
	Die Menge der Teilwörter von $w$ ist
	\begin{align*}
		\lrc{\lambda, a, b, c, ab, bc, abc},
	\end{align*}
	wobei nur $\lrc{\lambda, a, ab, abc}$ auch Präfixe von $w$ sind.
\end{myanswer}


\begin{myanswer}[mychallenge:0207]
	Wir dürfen annehmen, dass $x$ die Form $x = a_1a_2\ldots a_n$ hat. Wir werden die Anzahl der Teilwörter von $x$ der Länge $i = 1, i = 2$ bis zur Länge $i = n$ zählen.

	\underline{$i = 1$}:
	\begin{align*}
		1:\quad     & \blue{\underline{a_1}}a_2a_3a_4\ldots a_{n-3}a_{n-2}a_{n-1}a_n \\
		2:\quad     & a_1\blue{\underline{a_2}}a_3a_4\ldots a_{n-3}a_{n-2}a_{n-1}a_n \\
		3:\quad     & a_1a_2\blue{\underline{a_3}}a_4\ldots a_{n-3}a_{n-2}a_{n-1}a_n \\
		            & \ldots                                                         \\
		(n-3):\quad & a_1a_2a_3a_4\ldots \blue{\underline{a_{n-3}}}a_{n-2}a_{n-1}a_n \\
		(n-2):\quad & a_1a_2a_3a_4\ldots a_{n-3}\blue{\underline{a_{n-2}}}a_{n-1}a_n \\
		(n-1):\quad & a_1a_2a_3a_4\ldots a_{n-3}a_{n-2}\blue{\underline{a_{n-1}}}a_n \\
		n: \quad    & a_1a_2a_3a_4\ldots a_{n-3}a_{n-2}a_{n-1}\blue{\underline{a_n}}
	\end{align*}
	Offensichtlich gibt es $n$ viele Teilwörter der Länge $1$, nämlich genau die Teilwörter $\blue{a_1},\blue{a_2},\ldots ,\blue{a_n}$. Wir haben jeweils die Anfangsposition des Teilwortes in $x$ unterstrichen.

	\underline{$i = 2$}:
	\begin{align*}
		1: \quad    & \blue{\underline{a_1}a_2}a_3a_4\ldots a_{n-3}a_{n-2}a_{n-1}a_n \\
		2: \quad    & a_1\blue{\underline{a_2}a_3}a_4\ldots a_{n-3}a_{n-2}a_{n-1}a_n \\
		3: \quad    & a_1a_2\blue{\underline{a_3}a_4}\ldots a_{n-3}a_{n-2}a_{n-1}a_n \\
		            & \ldots                                                         \\
		(n-3):\quad & a_1a_2a_3a_4\ldots \blue{\underline{a_{n-3}}a_{n-2}}a_{n-1}a_n \\
		(n-2):\quad & a_1a_2a_3a_4\ldots a_{n-3}\blue{\underline{a_{n-2}}a_{n-1}}a_n \\
		(n-1):\quad & a_1a_2a_3a_4\ldots a_{n-3}a_{n-2}\blue{\underline{a_{n-1}}a_n}
	\end{align*}
	Es gibt $n-1$ viele Teilwörter der Länge $2$. Man beachte, dass sich die Anfangsposition der Teilwörter in $x$ um eine Position nach links verschoben hat, da das letzte Teilwort (ganz rechts) der Länge $2$ beim zweitletzten Symbol in $x$ beginnen muss.

	\underline{$i = 3$}:
	\begin{align*}
		1: \quad    & \blue{\underline{a_1}a_2a_3a_4}\ldots a_{n-3}a_{n-2}a_{n-1}a_n \\
		2: \quad    & a_1\blue{\underline{a_2}a_3a_4}\ldots a_{n-3}a_{n-2}a_{n-1}a_n \\
		            & \ldots                                                         \\
		(n-3):\quad & a_1a_2a_3a_4\ldots \blue{\underline{a_{n-3}}a_{n-2}a_{n-1}}a_n \\
		(n-2):\quad & a_1a_2a_3a_4\ldots a_{n-3}\blue{\underline{a_{n-2}}a_{n-1}a_n}
	\end{align*}
	Es gibt $n-2$ viele Teilwörter der Länge $3$.

	Nun nicht mehr schwierig zu sehen, dass es für ein Teilwort der Länge $i\in\lrc{1,2,\ldots ,n}$ genau $n-i+1$ mögliche Anfangspositionen in $x$ gibt, da die Positionen $n-i+2,3,\ldots ,n$ nicht möglich sind (\enquote{zu wenig Platz}). Wir summieren jetzt die Anzahl aller unterschiedlichen (nicht leeren) Teilwörter von $x$ der Längen $i=1,2,\ldots, n$:
	\begin{align*}
		\sum_{i=1}^n (n-i+1) = n(n+1)-\sum_{i=1}^n i = n(n+1) - \frac{n(n+1)}{2}=\frac{n(n+1)}{2}
	\end{align*}
	Hinzu kommt noch das leere Wort (welches Teilwort jedes Wortes ist) und als Resultat finden wir, dass es
	\begin{align*}
		\frac{n(n+1)}{2} + 1
	\end{align*}
	viele verschiedene Teilwörter von $x$ gibt.
\end{myanswer}

\begin{myanswer}[myexercise:0208]
	\begin{align*}
		L_3 = \lrc{a^4, a^2b, a^5, a^3b,ba^4, ba^2b}
	\end{align*}
	Bitte beachten Sie unbedingt, dass zwar $aa^4 = a^5$ aber $ba^2b \neq a^2b^2$, da $ba^2b = baab \neq a^2b^2 = aabb$.
\end{myanswer}

\begin{myanswer}[mychallenge:0209]
	Wir wählen $\Sigma = \lrc{0}$ und
	\begin{align*}
		L_1 = \lrc{0^1,0^2,\ldots ,0^k} = \setcm{0^i}{1\leq i \leq k}
	\end{align*}
	und
	\begin{align*}
		L_2 = \lrc{\lambda,0}.
	\end{align*}
	Dann gilt
	\begin{align*}
		L_1\cdot L_2 & = L_1\cdot \lrc{\lambda}\cup L_1\cdot \lrc{0} \\
		             & = L_1 \cup \setcm{0^i0}{1\leq i\leq k}        \\
		             & = L_1 \cup \setcm{0^i}{2\leq i \leq k+1}      \\
		             & = \setcm{0^i}{1\leq i\leq k+1},
	\end{align*}
	und somit $\abs{L_1L_2} = k+1$.

	Alternativ könnten Sie $L_2$ auch definieren als $L_2 = \lrc{a, aa}$ (ohne $L_1$ zu verändern).
\end{myanswer}

\clearpage
\shipoutAnswer


\clearpage
\section{Kapiteltest}

\begin{myexercise}\label{myexercise:kp00}
	Markieren Sie die zutreffende(n) Antwort(en).
	\begin{aenum}
		\item Es sei $\Sigma$ ein Alphabet.
		\begin{todolist}
			\item $\Sigma^*$ ist eine unendliche Menge.
			\item $\Sigma^*$ ist eine endliche Menge.
			\item $\Sigma^*$ ist keine Menge.
		\end{todolist}
		\item Es seien $\Sigma_1, \Sigma_2$ Alphabete. Seien $L_1$ eine endliche Sprache über $\Sigma_1$ und $L_2$ eine endliche Sprache über $\Sigma_2$. Dann gilt
		\begin{todolist}
			\item $\abs{L_1L_2} \leq \abs{L_1}\cdot \abs{L_2}$.
			\item $\abs{L_1L_2} = \abs{L_1}\cdot \abs{L_2}$.
			\item $\abs{L_1L_2} > \abs{L_1}\cdot \abs{L_2}$.
			\item $\abs{L_1L_2} \neq \abs{L_1}\cdot \abs{L_2}$.
		\end{todolist}
		\item
		\begin{todolist}
			\item $\lambda$ ist Element jedes Alphabets
			\item $\lambda$ ist ein Wort über jedem Alphabet
			\item $\abs{\lambda} = 0$
			\item $\lambda$ ist Suffix jedes Wortes über jedem Alphabet
		\end{todolist}
	\end{aenum}
\end{myexercise}


\begin{myexercise}\label{myexercise:kp01}
	Gegeben sei eine natürliche Zahl $n$. Bestimmen Sie die Anzahl der Wörter über $\Sigma_{\text{bool}}$ der Länge $n$, welche
	\begin{aenum}
		\item das Teilwort $01$ nicht enthalten.
		\item weder das Teilwort $01$ noch das Teilwort $00$ enthalten.
		\item das Teilwort $00$ nicht enthalten (schwierig!).
	\end{aenum}
\end{myexercise}

\begin{myexercise}\label{myexercise:kp02}
	\begin{aenum}
		\item Es sei $L := \lrc{a, ab, ba}$. Bestimmen Sie die Sprache $L^2$. Zur Erinnerung: $L^2 := L\cdot L$ (Konkatenation einer Sprache mit sich selbst).
		\item Existiert eine nichtleere endliche Sprache $L\neq\lrc{\lambda}$ über $\Sigma_{\text{bool}}$, welche die Gleichung
		\begin{align*}
			L = L^2
		\end{align*}
		erfüllt? Begründen Sie Ihre Antwort.
	\end{aenum}
\end{myexercise}

\begin{myexercise}\label{myexercise:kp03}
	Beweisen oder widerlegen Sie folgende Gleichung:
	\begin{align*}
		\lr{\lrc{0, 1}^2}^* = \lr{\lrc{0, 1}^*}^2.
	\end{align*}
\end{myexercise}

\clearpage
\section{Lösungen zum Kapiteltest}
\shipoutAnswer
\begin{myanswer}[myexercise:kp00]
	\begin{aenum}
		\item Es sei $\Sigma$ ein Alphabet.
		\begin{todolist}
			\item[\done] $\Sigma^*$ ist eine unendliche Menge.
			\item $\Sigma^*$ ist eine endliche Menge.
			\item $\Sigma^*$ ist keine Menge.
		\end{todolist}
		\item Es seien $\Sigma_1, \Sigma_2$ Alphabete. Seien $L_1$ eine endliche Sprache über $\Sigma_1$ und $L_2$ eine endliche Sprache über $\Sigma_2$. Dann gilt
		\begin{todolist}
			\item[\done] $\abs{L_1L_2} \leq \abs{L_1}\cdot \abs{L_2}$.
			\item $\abs{L_1L_2} = \abs{L_1}\cdot \abs{L_2}$.
			\item $\abs{L_1L_2} > \abs{L_1}\cdot \abs{L_2}$.
			\item $\abs{L_1L_2} \neq \abs{L_1}\cdot \abs{L_2}$.
		\end{todolist}
		\item
		\begin{todolist}
			\item $\lambda$ ist Element jedes Alphabets
			\item[\done] $\lambda$ ist ein Wort über jedem Alphabet
			\item[\done] $\abs{\lambda} = 0$
			\item[\done] $\lambda$ ist Suffix jedes Wortes über jedem Alphabet
		\end{todolist}
	\end{aenum}
\end{myanswer}

\begin{myanswer}[myexercise:kp01]
	\begin{aenum}
		\item\label{aenum:teila} Enthält ein binäres Wort $w$ der Länge $n\in\N$ nicht $01$ als Teilwort, dann hat $w$ die Form $w = 1^k0^m$ für zwei natürliche Zahlen $k, m$, welche die Gleichung $k + m = n$ erfüllen. Somit ist ein solches Wort durch die Wahl von $k$ eindeutig bestimmt. Für $k$ gelten die Ungleichungen $0\leq k \leq n$. Somit gibt es genau $n+1$ mögliche Wahlmöglichkeiten für $k$ und damit also genau $n+1$ solche (gesuchten) Wörter der Länge $n$.
		\item Wir haben in Teil \ref{aenum:teila} gesehen, dass jedes Wort mit Länge, welches $01$ nicht als Teilwort enthält, die Form $1^k0^m$ haben muss mit $k + m = n$. Alle Wörter für $n\leq 1$ erfüllen diese Bedingung. Es gibt ein Wort der Länge 0 und zwei Wörter der Länge 1, welche diese Bedingung erfüllen. Sei nun $n \geq 2$. Die einzigen solchen Wörter, welche nicht $00$ als Teilwort enthalten, sind $1^{n-1}0$ und $1^n$. Also, gibt es für $n\geq 2$ genau zwei solche Wörter.
	\end{aenum}
\end{myanswer}

\begin{myanswer}[myexercise:kp02]
	\begin{aenum}
		\item Es sei $L := \lrc{a, ab, ba}$. Dann ist $L^2$ gegeben durch
		\begin{align*}
			L^2 = \lrc{a^2, a^2b, aba, (ab)^2, ab^2a, ba^2, ba^2b, (ba)^2}.
		\end{align*}

		\item Nein, eine solche Sprache existiert nicht. Angenommen es gäbe eine solche Sprache $L$. Da $L$ nichtleer ist, existiert in $L$ ein Wort $w\in L$ maximaler Länge, also $\abs{w} = \max\setcm{\abs{u}}{u\in L}$ und es gilt $\abs{w}\geq 1$. Da $L^2 = L$, gilt $w^2\in L$. Doch dann gilt $\abs{w^2}= 2\abs{w}>\abs{w}$ und somit ist $w$ nicht das längste Wort in $L$. Dies ist ein Widerspruch.
	\end{aenum}
\end{myanswer}

\begin{myanswer}[myexercise:kp03]
	Die beiden Mengen sind verschieden. Offensichtlich gilt $0\in \lr{\lrc{0, 1}^*}^2$ aber $0\notin \lr{\lrc{0, 1}^2}^*$.
\end{myanswer}
